% JAC1964-20261121, /root/arithmetic_sign.  Direct application and limitation.
\subsection*{AR-22. What divisor Hodge index gives on the curve test}

\begin{proposition}[AR-22]
Let $W\subset A$ be a geometrically integral curve on a polarized
abelian threefold $(A,L)$ over a number field $K$, and assume
$L\operatorname{Ce}(W)=0$ in rational Chow groups.  Let
$f:C\to W\subset A$ be its normalization, put $X=C\times C$, and
equip $L$ with its canonical rigidified metric.  On $X$ set
\[
 \overline L_i=p_i^*f^*\overline L,\qquad
 \overline P=\overline L_1+\overline L_2,\qquad
 \overline Q=(f\times f)^*
 \bigl(\mu_+^*\overline L-p_1^*\overline L-p_2^*\overline L\bigr).
\]
Write $\langle\cdot\rangle_K=[K:\mathbf Q]^{-1}\widehat{\deg}(\cdot)$
for arithmetic intersections on $X$, and define
\[
 d=P^2>0,\qquad t=-\frac{Q^2}{P^2},\qquad
 B_0=\langle\overline P^3\rangle_K,\quad
 B_2=\langle\overline P\overline Q^2\rangle_K,\quad
 B_3=\langle\overline Q^3\rangle_K.
\]
Then $0\le t\le1$, $B_0\ge0$, and for every real $a\in[-1,1]$,
\[
 B_2+aB_3+a^2(2tB_2+t^2B_0)\le0.                 \tag{AR.22a}
\]
In particular,
\[
 B_2\le0,\qquad
 |B_3|\le -(1+2t)B_2-t^2B_0.                    \tag{AR.22b}
\]
These inequalities are symmetric under $B_3\mapsto-B_3$ and do not
orient the comparison of AR-18.  Precisely,
\[
 H_\pm(W)=B_0+3B_2\pm B_3,\qquad
 h_K(\operatorname{Ce}(W),\operatorname{Ce}(W))
       =-\frac{B_3}{6h^0(A,L)}.                  \tag{AR.22c}
\]
The desired semipositivity would require $B_3\le0$ for every such $W$.
\end{proposition}

\begin{proof}
The product $X$ is smooth, projective, and geometrically integral.
The line bundle $f^*L$ has positive degree on $C$, so $P$ is ample.
The Poincare bundle has degree zero on each curve fibre, whence
$PQ=0$ geometrically.  The ordinary Hodge index theorem on $X$ gives
$Q^2\le0$.  The theorem of the cube, with canonical metrics, gives
\[
 \overline R_\pm:=\overline P\pm\overline Q
       =(\mu_\pm\circ(f\times f))^*\overline L.
\]
These are nef adelic line bundles; so is $\overline P$.  Since
$R_\pm^2=P^2+Q^2\ge0$, we obtain $0\le t\le1$.
Nonnegativity of the top intersection of nef adelic bundles gives
$B_0\ge0$.

The primitive metric cancellation proved in AR-18 applies to the
normalization and gives the exact arithmetic equality
\[
 \langle\overline P^2\overline Q\rangle_K=0.       \tag{AR.22d}
\]
For clarity, this is more than the geometric equality $PQ=0$: it
uses the triviality, in the rational adelic Picard group, of the
relative Deligne pairing at the Albanese sum
$\operatorname{alb}_A(L\cdot W)=0$.  We use precisely that already
proved cancellation, not a numerical replacement for primitivity.

For a rational $a$ with $|a|<1$, put
\[
 \overline R_a=\overline P+a\overline Q,
 \qquad \overline M_a=\overline Q+at\overline P.
\]
The metric on $\overline R_a$ is nef, because
\[
 \overline R_a=\frac{1+a}{2}\overline R_+
                    +\frac{1-a}{2}\overline R_-.
\]
Its underlying line bundle is ample: for $a\ne0$ it is
$(1-|a|)P+|a|R_{\operatorname{sgn}(a)}$, the sum of a positive
multiple of an ample class and a nef class; the case $a=0$ is $P$.
In particular the underlying $R_a$ is big.  The line bundle
$\overline Q=\overline R_+-\overline P$ is integrable, and hence
so is $\overline M_a$.  Finally the geometric intersection is
\[
 M_aR_a=(Q+atP)(P+aQ)=a(Q^2+tP^2)=0.
\]

Apply Yuan--Zhang's adelic divisor Hodge index theorem, published
Theorem 3.2 in \emph{Math. Ann.} \textbf{367} (2017), pp.~1134--1135,
to $\overline M_a$ and the one nef bundle $\overline R_a$ on the
normal integral surface $X_{\overline K}$.  Its inequality requires
integrability of $\overline M_a$, nefness of $\overline R_a$, bigness
of its underlying line bundle, and $M_aR_a=0$, all just checked.
It yields
$\langle\overline M_a^2\overline R_a\rangle_K\le0$.
The extra arithmetic positivity and boundedness hypotheses in the
theorem concern its equality characterization and are not needed
here; no equality characterization is invoked.  By (AR.22d),
\[
 \begin{split}
 \langle(\overline Q+at\overline P)^2
                     (\overline P+a\overline Q)\rangle_K
 &=B_2+aB_3+2a^2tB_2+a^2t^2B_0.
 \end{split}
\]
This proves (AR.22a) for rational $|a|<1$.  Polynomial continuity
gives it on the closed real interval; this limit avoids any need to
assume $R_+$ or $R_-$ itself is big.  Setting $a=0,1,-1$ gives
(AR.22b).  Expanding $\langle(\overline P\pm\overline Q)^3\rangle_K$
and using proper pushforward and (AR.22d) proves the first formula
in (AR.22c); AR-18 proves the second.
\end{proof}

\paragraph{The exact extent of this test.}
The whole family (AR.22a), not just its two endpoints, survives
unchanged under $B_3\mapsto-B_3$ by replacing $a$ with $-a$.
The separate nonnegativities $H_\pm(W)\ge0$ also survive that
replacement.  These numerical consequences therefore provide size
bounds but no preferred sign of $B_3$.  For example, the formal
values $t=1/3$, $B_0=4$, $B_2=-1$, $B_3=\pm1/10$ satisfy the
entire displayed inequality family and both inequalities
$H_\pm\ge0$.  This is only a demonstration of the logical symmetry
of these constraints, not an assertion that these numbers are
realized by a curve.

For an arbitrary primitive odd class $z$, AR-21 supplies a curve
$W$ and $M>0$ with $\operatorname{Ce}(W)=Mz$.  The resulting bound
is
\[
 |h_K(z,z)|\le
 \frac{-(1+2t)B_2-t^2B_0}{6h^0(A,L)M^2}.          \tag{AR.22e}
\]
The auxiliary quantities on the right depend on the chosen curve
representative.  Nothing here asserts that they are invariant under
rational equivalence, or that bounding them controls their sign.
The full comparison $H_-(W)\ge H_+(W)$ remains the next gate.

\paragraph{Primary locator.}
The complete published source was retrieved from the author's site,
\url{https://web.math.princeton.edu/\string~shouwu/publications/hodge-index}.
The relevant statement is Theorem 3.2, with conventions in
\S3.1 and the appendix on arithmetic intersections.  The earlier
May 21, 2013 manuscript calls the corresponding result Theorem 1.3;
the numbering and page locators used above refer to the published
2017 paper.  Only its inequality is used.  The geometric index
theorem and the canonical metric identities are background; the
choice of $M_a,R_a$ and the symmetric bounds are the deductions here.
